% A revista exige resumo e palavras-chave em português e em uma língua estrangeira.
% Não há limite de palavras especificado nas diretrizes gerais consultadas.
% Os textos abaixo são exemplos de apresentação, não resumos finais da pesquisa.
\section*{Resumo}
Em forma de diário e carta de intenções, este ensaio propõe pensar o \enquote{Antropoceno} como uma crise da comunicação e da constituição de comunalidades, articulando a cosmofobia, a Constituição moderna e os usos das aspas em diferentes regimes de conhecimento. Em diálogo com Manuela Carneiro da Cunha, Nego Bispo, Ailton Krenak e Bruno Latour, apresenta Cecinepá como uma heterotopia especulativa na qual saber, afeto e laço compõem uma cartografia das diferenças entre um corpo e sua participação em outros corpos. A própria escrita e seus rodapés são propostos como uma experiência de práticas científicas orientadas pela e para a comunicação, voltadas à composição de um mundo comum entre diferentes modos de conhecer e existir.

\noindent\textbf{Palavras-chave:} Antropoceno; comunicação; Cecinepá.

\section*{Abstract}
\begin{otherlanguage}{english}
This template organizes an essay as a sequence of diary entries in which a statement of research intentions concerning the ``Anthropocene'' is developed. Footnotes provide space to carry out the procedures and develop the references and arguments announced in the main text.

\noindent\textbf{Keywords:} Anthropocene; communication; naming; Cecinepá.
\end{otherlanguage}
